\chapter{共轭理论与无穷维对偶}

\label{chap:06}

第 \ref{chap:05} 章把下半连续性、可微性与次微分组织成了凸泛函的局部语言，但这些对象仍然主要生活在原空间 $X$ 中。真正的对偶理论则要求把问题整体搬到对偶空间里重新观察：一个泛函如何通过支撑超平面被编码为另一个泛函？原问题的信息在对偶空间中会损失多少？何时对偶问题只给出下界，何时又能完整恢复原问题的最优值与最优性条件？这些问题的统一答案，正是本章要建立的共轭理论框架。

本章沿着一条清晰的链条推进。第 \ref{sec:6-1} 节首先定义 Fenchel 共轭，并证明 Fenchel--Young 不等式与其次微分取等条件，说明“共轭”本质上是在对偶空间中重新记录支撑信息。第 \ref{sec:6-2} 节继续引入双共轭，证明 Fenchel--Moreau 定理，并解释为什么对偶法看到的往往是原问题的闭凸正则化版本。第 \ref{sec:6-3} 节转向下确界卷积，给出它与共轭、Moreau 包络和变量分裂之间的代数联系，为第 9 至 10 章算法准备接口。最后第 \ref{sec:6-4} 节把这些工具收束成 Fenchel--Rockafellar 对偶框架，用统一的原--对偶模板、资格条件和对偶间隙语言，为第 \ref{chap:07} 章 KKT 条件、第 \ref{chap:08} 章有限维坍缩以及后续应用章节提供总入口。

阅读这一章时，需要始终把第 \ref{sec:4-3} 节的分离定理与第 \ref{sec:5-3} 节的次微分语言同时放在脑中。没有 Hahn--Banach 分离，便没有支撑超平面的丰富性；没有次微分，便无法把共轭取等条件重写为最优性条件。另一方面，第 \ref{sec:5-4} 节 Moreau--Rockafellar 和公式已经预示：资格条件不是证明中的装饰，而是决定“原空间与对偶空间能否无损交换信息”的真正门槛。后文关于广义 KKT、极大单调算子、ADMM 与几何规划的许多关键步骤，都会回链到这一章。

\section{Fenchel 共轭与 Fenchel-Young 不等式}

\label{sec:6-1}

共轭理论的出发点是一个简单但极其深刻的观察：若一个凸泛函由其所有支撑超平面决定，那么与其直接研究函数值本身，不如把这些支撑超平面收集起来，看成活在对偶空间中的新对象。这样得到的对象就是 Fenchel 共轭。它既是第 \ref{sec:4-3} 节分离理论的函数化版本，又是第 \ref{sec:5-3} 节次微分理论的对偶表达，因此本节实际上承担着把“支撑超平面语言”升级为“对偶泛函语言”的任务。

\subsection{Fenchel 共轭的定义}

\begin{definition}[Fenchel 共轭]\label{def:fenchel-conjugate}
设 $X$ 为 Banach 空间，$f\colon X\to \mathbb R\cup\{+\infty\}$ 为 proper 泛函。其\emph{Fenchel 共轭}定义为函数
\[
f^\ast\colon X^\ast\to \mathbb R\cup\{+\infty\},
\qquad
f^\ast(x^\ast):=\sup_{x\in X}\{\langle x^\ast,x\rangle-f(x)\}.
\]
这里的上确界允许取值 $+\infty$。若 $f$ 为凸泛函，则 $f^\ast$ 也是凸且下半连续的。
\end{definition}

\begin{remark}
定义~\ref{def:fenchel-conjugate} 说明，共轭并不是把原函数“再写一遍”，而是把所有可能的仿射下界
\[
x\mapsto \langle x^\ast,x\rangle-c
\]
压缩成对偶变量 $x^\ast$ 的函数。换言之，$f^\ast(x^\ast)$ 衡量的是斜率为 $x^\ast$ 的仿射函数要支撑 $f$，其截距至少需要抬高多少。
\end{remark}

\subsection{Fenchel-Young 不等式与取等条件}

\begin{proposition}[Fenchel--Young 不等式]\label{prop:fenchel-young-inequality}
设 $f\colon X\to \mathbb R\cup\{+\infty\}$ 为 proper 泛函。则对任意 $x\in X$ 与任意 $x^\ast\in X^\ast$，都有
\[
f(x)+f^\ast(x^\ast)\ge \langle x^\ast,x\rangle.
\]
\end{proposition}

\begin{proof}
由定义~\ref{def:fenchel-conjugate}，
\[
f^\ast(x^\ast)=\sup_{y\in X}\{\langle x^\ast,y\rangle-f(y)\}.
\]
特别地，把 $y$ 取为当前的 $x$，便得到
\[
f^\ast(x^\ast)\ge \langle x^\ast,x\rangle-f(x).
\]
移项即得
\[
f(x)+f^\ast(x^\ast)\ge \langle x^\ast,x\rangle.
\]
\end{proof}

\begin{proposition}[Fenchel--Young 取等与次微分]\label{prop:fenchel-young-equality-subgradient}
设 $f\colon X\to \mathbb R\cup\{+\infty\}$ 为 proper 凸泛函，$x\in \operatorname{dom}f$，$x^\ast\in X^\ast$。则下列命题等价：
\begin{enumerate}
  \item $x^\ast\in \partial f(x)$；
  \item
  \[
  f(x)+f^\ast(x^\ast)=\langle x^\ast,x\rangle;
  \]
  \item $x\in \partial f^\ast(x^\ast)$。
\end{enumerate}
\end{proposition}

\begin{proof}
由定义~\ref{def:subdifferential}，$x^\ast\in \partial f(x)$ 当且仅当对任意 $y\in X$，
\[
f(y)\ge f(x)+\langle x^\ast,y-x\rangle.
\]
移项后得
\[
\langle x^\ast,y\rangle-f(y)\le \langle x^\ast,x\rangle-f(x),\qquad \forall y\in X.
\]
对左边取上确界，得到
\[
f^\ast(x^\ast)\le \langle x^\ast,x\rangle-f(x).
\]
再结合命题~\ref{prop:fenchel-young-inequality}，即得
\[
f^\ast(x^\ast)=\langle x^\ast,x\rangle-f(x),
\]
即 (1)$\Rightarrow$(2)。

反过来，若 (2) 成立，则对任意 $y\in X$，
\[
\langle x^\ast,y\rangle-f(y)\le f^\ast(x^\ast)=\langle x^\ast,x\rangle-f(x),
\]
从而
\[
f(y)\ge f(x)+\langle x^\ast,y-x\rangle.
\]
故 $x^\ast\in \partial f(x)$。这给出 (2)$\Rightarrow$(1)。

把上一段论证应用于对偶空间上的凸泛函 $f^\ast$，并注意共轭再取一次所得到的变量互换，便得到 (2) 与 (3) 的等价性。
\end{proof}

\begin{remark}
命题~\ref{prop:fenchel-young-equality-subgradient} 是本章最重要的“翻译器”之一：它把原空间里的次微分条件 $x^\ast\in\partial f(x)$ 翻译成共轭表达式的取等条件。后文第 \ref{sec:6-4} 节的对偶定理和第 \ref{sec:7-4} 节的 KKT 条件，本质上都在重复使用这一翻译。
\end{remark}

\subsection{典型例子}

\begin{example}[指示函数与支撑函数]\label{ex:indicator-support-conjugate}
设 $C\subset X$ 为非空集合，指示函数定义为
\[
\iota_C(x)=
\begin{cases}
0, & x\in C,\\
+\infty, & x\notin C.
\end{cases}
\]
则其共轭为
\[
\iota_C^\ast(x^\ast)=\sup_{x\in C}\langle x^\ast,x\rangle=: \sigma_C(x^\ast),
\]
这里 $\sigma_C$ 称为 $C$ 的支撑函数。把这个例子放在这里，是为了把第 \ref{sec:4-3} 节的支撑超平面几何直接转写为函数等式：集合的约束信息，在对偶空间中正是由支撑函数承载。
\end{example}

\begin{example}[范数与对偶范数的共轭]\label{ex:norm-dual-norm-conjugate}
设 $\|\cdot\|$ 是 $X$ 上的范数，$\|\cdot\|_\ast$ 是对偶范数。对任意 $\lambda>0$，令
\[
f(x)=\lambda \|x\|.
\]
则
\[
f^\ast(x^\ast)=
\begin{cases}
0, & \|x^\ast\|_\ast\le \lambda,\\
+\infty, & \|x^\ast\|_\ast>\lambda.
\end{cases}
\]
把这个例子放在这里，是为了说明范数正则化在对偶空间中对应的是一个对偶球约束；这正是 Boyd 中“范数球与对偶球”图像的无穷维上位版本。
\end{example}

\begin{proof}
对任意 $x^\ast\in X^\ast$，
\[
f^\ast(x^\ast)=\sup_{x\in X}\{\langle x^\ast,x\rangle-\lambda\|x\|\}.
\]
若 $\|x^\ast\|_\ast\le \lambda$，则由对偶范数定义
\[
\langle x^\ast,x\rangle-\lambda\|x\|
\le (\|x^\ast\|_\ast-\lambda)\|x\|\le 0,
\]
且在 $x=0$ 处取到 $0$，故 $f^\ast(x^\ast)=0$。反之若 $\|x^\ast\|_\ast>\lambda$，则由 Hahn--Banach 支撑论证可取单位向量方向 $u$ 使 $\langle x^\ast,u\rangle>\lambda$，于是
\[
\langle x^\ast,tu\rangle-\lambda\|tu\|
=t(\langle x^\ast,u\rangle-\lambda)\to +\infty
\qquad (t\to+\infty).
\]
故 $f^\ast(x^\ast)=+\infty$。
\end{proof}

\begin{example}[正则化项的对偶表示]
在 Hilbert 空间 $H$ 上考虑
\[
f(u)=\frac{1}{2}\|u\|^2.
\]
由第 \ref{sec:3-1} 节定理~\ref{thm:riesz-representation-hilbert} 可把 $H^\ast$ 与 $H$ 识别，于是
\[
f^\ast(v)=\frac{1}{2}\|v\|^2.
\]
把这个例子放在这里，是为了给后文 Moreau 包络、prox 映射和二次正则化提供最核心的对偶基线。
\end{example}

\subsection*{有限维对照}

Boyd 书中关于共轭函数的几何图像，主要依赖于二维或三维里“用斜率看支撑线”的直觉。本节说明，这个直觉在无穷维里依然成立，但真正稳固的形式是定义~\ref{def:fenchel-conjugate} 与命题~\ref{prop:fenchel-young-equality-subgradient}。特别是例~\ref{ex:indicator-support-conjugate} 和例~\ref{ex:norm-dual-norm-conjugate} 表明，Boyd 里熟悉的支撑函数、对偶球和范数正则化，都只是同一套对偶编码机制在有限维 Hilbert 空间中的坍缩。

\subsection*{习题（待补）}

\begin{exercise}[Fenchel--Young 取等占位]\label{exr:fenchel-young-placeholder}
补一题：对一个具体的凸泛函验证命题~\ref{prop:fenchel-young-equality-subgradient}，并把取等条件写成第 \ref{sec:5-3} 节次微分语言中的最优性判断。
\end{exercise}

\begin{exercise}[支撑函数桥接题占位]\label{exr:support-function-placeholder}
补一题：从例~\ref{ex:indicator-support-conjugate} 出发，计算一个有限维凸体的支撑函数，并说明该结果如何与 Boyd 书中的超平面支撑图像对应。
\end{exercise}


\section{双共轭与闭凸包}

\label{sec:6-2}

一旦有了共轭，最自然的问题就是：再共轭一次会发生什么？有限维图像直觉会让人猜测“应该回到原函数”，但无穷维对偶理论告诉我们，真正回来的并不是原问题的全部细节，而是它的闭凸正则化版本。双共轭正是刻画这种“信息损失最小的闭凸修正”的工具。因此，本节的核心不仅是证明 Fenchel--Moreau 定理（下文也写作 Fenchel-Moreau 定理），更是解释为什么对偶法在某些时候只能看到松弛后的问题，而看不到原始的非闭、非凸细节。第 \ref{sec:5-1} 节定义~\ref{def:lower-semicontinuity-functional} 中的下半连续性，在这里正被重新识别为“函数能够被双共轭完整恢复”的必要正则性。

\subsection{双共轭的定义与基本不等式}

\begin{definition}[双共轭]\label{def:biconjugate-functional}
设 $f\colon X\to \mathbb R\cup\{+\infty\}$ 为 proper 泛函。定义其\emph{双共轭}为
\[
f^{\ast\ast}(x):=\sup_{x^\ast\in X^\ast}\{\langle x^\ast,x\rangle-f^\ast(x^\ast)\},
\qquad x\in X.
\]
也就是说，$f^{\ast\ast}$ 是把 $f^\ast$ 看成定义在对偶空间上的泛函后，再次对其取 Fenchel 共轭所得到的对象。
\end{definition}

\begin{proposition}[双共轭总在原函数之下]\label{prop:biconjugate-below-function}
设 $f\colon X\to \mathbb R\cup\{+\infty\}$ 为 proper 泛函。则对任意 $x\in X$，
\[
f^{\ast\ast}(x)\le f(x).
\]
\end{proposition}

\begin{proof}
由命题~\ref{prop:fenchel-young-inequality}，对任意 $x\in X$ 与任意 $x^\ast\in X^\ast$，
\[
f(x)+f^\ast(x^\ast)\ge \langle x^\ast,x\rangle.
\]
移项得
\[
\langle x^\ast,x\rangle-f^\ast(x^\ast)\le f(x).
\]
对左边关于 $x^\ast$ 取上确界，便得到
\[
f^{\ast\ast}(x)\le f(x).
\]
\end{proof}

\subsection{Fenchel--Moreau 定理}

\begin{theorem}[Fenchel--Moreau 定理]\label{thm:fenchel-moreau}
设 $X$ 为 Banach 空间，$f\colon X\to \mathbb R\cup\{+\infty\}$ 为 proper、凸且下半连续的泛函。则
\[
f^{\ast\ast}=f.
\]
\end{theorem}

\begin{proof}
由命题~\ref{prop:biconjugate-below-function}，总有
\[
f^{\ast\ast}\le f.
\]
只需证明反向不等式。固定 $x\in X$ 与实数 $r<f(x)$。考虑乘积空间 $X\times \mathbb R$ 中的点 $(x,r)$ 与上图像
\[
\operatorname{epi}f:=\{(u,t)\in X\times \mathbb R:t\ge f(u)\}.
\]
由于 $f$ 凸且下半连续，$\operatorname{epi}f$ 是闭凸集，而 $(x,r)\notin \operatorname{epi}f$。由第 \ref{sec:4-3} 节定理~\ref{thm:strict-separation}，存在非零连续线性泛函 $(p,\lambda)\in X^\ast\times \mathbb R$ 与常数 $c$，使
\[
\langle p,u\rangle+\lambda t\ge c>\langle p,x\rangle+\lambda r,
\qquad \forall (u,t)\in \operatorname{epi}f.
\]
由于若 $(u,t)\in \operatorname{epi}f$ 则 $(u,t+s)\in \operatorname{epi}f$ 对所有 $s\ge 0$ 仍成立，故不可能有 $\lambda<0$；否则令 $s\to+\infty$ 将破坏左侧不等式。若 $\lambda=0$，则
\[
\langle p,u\rangle\ge c>\langle p,x\rangle
\]
对所有 $u\in \operatorname{dom}f$ 成立，这与 $f$ 的凸下半连续结构结合时会把 $(x,r)$ 完全隔离到定义域外；在当前取定的 $r<f(x)$ 情形下，我们可直接把 $(x,f(x))\in\operatorname{epi}f$ 代入，得到
\[
\langle p,x\rangle+\lambda f(x)\ge c>\langle p,x\rangle+\lambda r,
\]
从而 $\lambda>0$。

于是对任意 $u\in X$，令 $t=f(u)$，有
\[
\langle p,u\rangle+\lambda f(u)\ge c>\langle p,x\rangle+\lambda r.
\]
移项并除以 $\lambda$ 得
\[
r<f(u)+\left\langle \frac{p}{\lambda},u-x\right\rangle,\qquad \forall u\in X.
\]
令
\[
x^\ast:=-\frac{p}{\lambda},
\]
则
\[
r< f(u)-\langle x^\ast,u\rangle+\langle x^\ast,x\rangle,\qquad \forall u\in X.
\]
对右边关于 $u$ 取下确界，便得到
\[
r\le \langle x^\ast,x\rangle-f^\ast(x^\ast)\le f^{\ast\ast}(x).
\]
由于这对每个 $r<f(x)$ 都成立，令 $r\uparrow f(x)$，便得
\[
f(x)\le f^{\ast\ast}(x).
\]
综上
\[
f(x)=f^{\ast\ast}(x),\qquad \forall x\in X.
\]
\end{proof}

\begin{remark}
定理~\ref{thm:fenchel-moreau} 的真正含义，不是“双共轭把函数又变回来”这样一句形式口号，而是：当且仅当原函数已经是 proper、闭且凸的，它才会完整地被对偶支撑信息重建。换言之，对偶法看到的是闭凸正则化后的问题；原问题若不闭或不凸，双共轭只会保留其对偶侧可见的那部分结构。
\end{remark}

\subsection{上图像闭凸包与闭凸正则化}

\begin{proposition}[双共轭的上图像闭凸包表述]\label{prop:epigraph-biconjugate-closure}
设 $f\colon X\to \mathbb R\cup\{+\infty\}$ 为 proper 泛函。则 $f^{\ast\ast}$ 是不超过 $f$ 的最大下半连续凸泛函；等价地，
\[
\operatorname{epi}(f^{\ast\ast})=\overline{\operatorname{conv}(\operatorname{epi}f)},
\]
其中闭包取于 $X\times \mathbb R$ 的乘积拓扑。
\end{proposition}

\begin{proof}
命题~\ref{prop:biconjugate-below-function} 已说明 $f^{\ast\ast}\le f$，而由定义~\ref{def:biconjugate-functional} 可知 $f^{\ast\ast}$ 总是凸且下半连续。若 $g\le f$ 且 $g$ 为下半连续凸泛函，则由共轭定义可得
\[
g^\ast\ge f^\ast.
\]
再取一次共轭，便有
\[
g=g^{\ast\ast}\le f^{\ast\ast},
\]
其中等号用到了定理~\ref{thm:fenchel-moreau} 对 $g$ 的应用。因此 $f^{\ast\ast}$ 的确是最大的下半连续凸 minorant。将这一叙述翻译到上图像层面，正得到
\[
\operatorname{epi}(f^{\ast\ast})=\overline{\operatorname{conv}(\operatorname{epi}f)}.
\]
\end{proof}

\begin{corollary}[指示函数、支撑函数与闭凸包]\label{cor:indicator-support-bipolar}
设 $C\subset X$ 为非空集合，则
\[
(\iota_C)^{\ast\ast}=\iota_{\overline{\operatorname{conv}}C}.
\]
特别地，若 $C$ 已经是非空闭凸集，则由例~\ref{ex:indicator-support-conjugate}，
\[
\sigma_C^\ast=\iota_C.
\]
\end{corollary}

\begin{proof}
由例~\ref{ex:indicator-support-conjugate}，$\iota_C^\ast=\sigma_C$。另一方面，$\iota_C$ 的最大下半连续凸 minorant 正是 $\iota_{\overline{\operatorname{conv}}C}$，故由命题~\ref{prop:epigraph-biconjugate-closure} 得
\[
(\iota_C)^{\ast\ast}=\iota_{\overline{\operatorname{conv}}C}.
\]
若 $C$ 本身已经闭凸，则右侧就是 $\iota_C$，从而 $\sigma_C^\ast=\iota_C$。
\end{proof}

\subsection{典型例子}

\begin{example}[非闭凸函数的双共轭正则化]
设 $X=\mathbb R$，并定义
\[
f(x)=
\begin{cases}
0, & x>0,\\
1, & x=0,\\
+\infty, & x<0.
\end{cases}
\]
则 $f$ 是凸但在 $0$ 处并不下半连续，其双共轭满足
\[
f^{\ast\ast}(x)=
\begin{cases}
0, & x\ge 0,\\
+\infty, & x<0.
\end{cases}
\]
把这个例子放在这里，是为了说明双共轭不是“原函数复制器”，而是闭凸修补器：它自动把孤立的非闭性抹平。
\end{example}

\begin{example}[有限维上图像闭凸包]
在 $\mathbb R^n$ 中，命题~\ref{prop:epigraph-biconjugate-closure} 可以直接理解为“对函数上图像先取凸包再闭包”。把这个例子放在这里，是为了把 Boyd 中关于上图像和支撑超平面的几何直觉直接接回到双共轭公式；有限维图像之所以好画，只是因为闭凸包在那里可以被可视化，而不是因为结论只属于有限维。
\end{example}

\begin{example}[替代损失与闭凸正则化]
在机器学习中，0-1 损失的闭凸包会导向铰链损失、logistic 损失等 surrogate 目标。把这个例子放在这里，是为了说明“对偶法看到的是闭凸正则化后的问题”并非纯理论现象，而是现代统计学习中每天都在发生的建模动作。
\end{example}

\subsection*{有限维对照}

Boyd 书中常把闭凸包与对偶下界放在几何图像里解释，例如通过支撑超平面观察一个非凸对象被其凸包替代。本节把这种几何说法提升为定理~\ref{thm:fenchel-moreau} 与命题~\ref{prop:epigraph-biconjugate-closure}。在有限维里，这些结果看上去像是图像事实；在无穷维里，它们真正依赖的是下半连续性、分离定理和对偶支撑信息的完整性。

\subsection*{习题（待补）}

\begin{exercise}[双共轭桥接题占位]\label{exr:biconjugate-placeholder}
补一题：对一个不是下半连续的凸函数显式计算其双共轭，并说明双共轭在哪个点上修补了原函数的非闭性。
\end{exercise}

\begin{exercise}[闭凸包占位]\label{exr:epigraph-closure-placeholder}
补一题：从命题~\ref{prop:epigraph-biconjugate-closure} 出发，在有限维中把一个具体函数的上图像闭凸包画出来，并与其双共轭函数逐点对应。
\end{exercise}


\section{下确界卷积}

\label{sec:6-3}

前两节说明，共轭把“求和难题”部分转化成了对偶空间里的代数运算。但在建模和算法中，我们经常还想做另一件事：把两个泛函组合成一个更平滑或更易分裂的对象。下确界卷积正是这个目的的标准工具。它一方面把两个代价通过“变量拆分后再求最优拼接”的方式组合起来，另一方面又在共轭侧转化为简单的求和，因此既是 Moreau 包络与平滑化的核心来源，也是第 9 至 10 章算法中变量分裂思想的代数前身。

\subsection{下确界卷积的定义}

\begin{definition}[下确界卷积]\label{def:infimal-convolution}
设 $X$ 为 Banach 空间，$f,g\colon X\to \mathbb R\cup\{+\infty\}$ 为 proper 泛函。定义它们的\emph{下确界卷积}为
\[
(f\square g)(x):=\inf_{y\in X}\{f(y)+g(x-y)\},\qquad x\in X.
\]
这一运算也常记作 infimal convolution。
\end{definition}

\begin{proposition}[下确界卷积保持凸性]\label{prop:infimal-convolution-convexity}
若 $f,g$ 为 proper 凸泛函，则 $f\square g$ 仍是凸泛函。若再假设其中一个函数连续并且另一个下半连续，则 $f\square g$ 在相应点上也是下半连续的。
\end{proposition}

\begin{proof}
任取 $x_1,x_2\in X$ 与 $\lambda\in[0,1]$。对任意 $\varepsilon>0$，可选 $y_1,y_2\in X$ 使
\[
f(y_i)+g(x_i-y_i)\le (f\square g)(x_i)+\varepsilon,\qquad i=1,2.
\]
由凸性，
\[
f(\lambda y_1+(1-\lambda)y_2)\le \lambda f(y_1)+(1-\lambda)f(y_2),
\]
且
\[
g(\lambda(x_1-y_1)+(1-\lambda)(x_2-y_2))
\le
\lambda g(x_1-y_1)+(1-\lambda)g(x_2-y_2).
\]
注意
\[
\lambda(x_1-y_1)+(1-\lambda)(x_2-y_2)
=
(\lambda x_1+(1-\lambda)x_2)-(\lambda y_1+(1-\lambda)y_2).
\]
故
\[
(f\square g)(\lambda x_1+(1-\lambda)x_2)
\le
\lambda(f\square g)(x_1)+(1-\lambda)(f\square g)(x_2)+\varepsilon.
\]
令 $\varepsilon\downarrow 0$ 得凸性。下半连续性部分来自定义中的下确界与第 \ref{sec:5-4} 节定理~\ref{thm:moreau-rockafellar-sum} 相同的资格条件机制，这里不再展开一般证明。
\end{proof}

\subsection{共轭与下确界卷积交换}

\begin{theorem}[下确界卷积的共轭交换公式]\label{thm:conjugate-infimal-convolution}
设 $f,g\colon X\to \mathbb R\cup\{+\infty\}$ 为 proper 泛函，则对任意 $x^\ast\in X^\ast$，
\[
(f\square g)^\ast(x^\ast)=f^\ast(x^\ast)+g^\ast(x^\ast).
\]
\end{theorem}

\begin{proof}
由定义~\ref{def:infimal-convolution}，
\[
(f\square g)^\ast(x^\ast)
=
\sup_{x\in X}\left\{\langle x^\ast,x\rangle-\inf_{y\in X}[f(y)+g(x-y)]\right\}.
\]
交换上确界与下确界后，可写为
\[
(f\square g)^\ast(x^\ast)
=
\sup_{x,y\in X}\{\langle x^\ast,x\rangle-f(y)-g(x-y)\}.
\]
令 $z=x-y$，则 $x=y+z$，因此
\[
\langle x^\ast,x\rangle-f(y)-g(x-y)
=
\langle x^\ast,y\rangle-f(y)+\langle x^\ast,z\rangle-g(z).
\]
于是
\[
(f\square g)^\ast(x^\ast)
=
\sup_{y\in X}\{\langle x^\ast,y\rangle-f(y)\}
\;+\;
\sup_{z\in X}\{\langle x^\ast,z\rangle-g(z)\}.
\]
右边正是
\[
f^\ast(x^\ast)+g^\ast(x^\ast).
\]
\end{proof}

\begin{remark}
定理~\ref{thm:conjugate-infimal-convolution} 是本节的代数核心。它说明：原空间里的“先拆变量、再取下确界”在对偶空间里恰好变成简单的求和。第 \ref{sec:5-4} 节命题~\ref{prop:subdifferential-linear-pullback} 与定理~\ref{thm:moreau-rockafellar-sum} 告诉我们，次微分演算同样依赖这种“在原空间难、在对偶空间易”的交换机制。
\end{remark}

\subsection{Moreau 包络}

\begin{definition}[Moreau 包络]\label{def:moreau-envelope}
设 $H$ 为 Hilbert 空间，$f\colon H\to \mathbb R\cup\{+\infty\}$ 为 proper 泛函，$\lambda>0$。定义
\[
e_\lambda f(x):=\inf_{y\in H}\left\{f(y)+\frac{1}{2\lambda}\|x-y\|^2\right\}.
\]
该函数称为 $f$ 的\emph{Moreau 包络}。它正是 $f$ 与二次惩罚项
\[
q_\lambda(x):=\frac{1}{2\lambda}\|x\|^2
\]
的下确界卷积。
\end{definition}

\begin{proposition}[Moreau 包络的共轭]\label{prop:moreau-envelope-conjugate}
设 $H$ 为 Hilbert 空间，$f\colon H\to \mathbb R\cup\{+\infty\}$ 为 proper 泛函。则
\[
(e_\lambda f)^\ast(v)=f^\ast(v)+\frac{\lambda}{2}\|v\|^2,\qquad v\in H.
\]
\end{proposition}

\begin{proof}
由定义~\ref{def:moreau-envelope}，
\[
e_\lambda f=f\square q_\lambda.
\]
再由定理~\ref{thm:conjugate-infimal-convolution}，
\[
(e_\lambda f)^\ast=f^\ast+q_\lambda^\ast.
\]
因此只需计算 $q_\lambda^\ast$。对任意 $v\in H$，
\[
q_\lambda^\ast(v)
=
\sup_{x\in H}\left\{\langle v,x\rangle-\frac{1}{2\lambda}\|x\|^2\right\}.
\]
把右边配方可得
\[
\langle v,x\rangle-\frac{1}{2\lambda}\|x\|^2
=
-\frac{1}{2\lambda}\|x-\lambda v\|^2+\frac{\lambda}{2}\|v\|^2.
\]
故上确界在 $x=\lambda v$ 处取到，且
\[
q_\lambda^\ast(v)=\frac{\lambda}{2}\|v\|^2.
\]
代回即得
\[
(e_\lambda f)^\ast(v)=f^\ast(v)+\frac{\lambda}{2}\|v\|^2.
\]
\end{proof}

\subsection{典型例子}

\begin{example}[两个指标函数与 Minkowski 和]
设 $A,B\subset X$ 为非空集合。则
\[
(\iota_A\square \iota_B)(x)=
\begin{cases}
0, & x\in A+B,\\
+\infty, & x\notin A+B,
\end{cases}
\]
其中
\[
A+B:=\{a+b:a\in A,\ b\in B\}
\]
是 Minkowski 和。把这个例子放在这里，是为了说明下确界卷积是集合运算的函数化版本：在指标函数层面，它正对应于集合的加法。
\end{example}

\begin{example}[Huber 平滑与 Moreau 包络]
在 $\mathbb R$ 上取 $f(t)=|t|$，则其 Moreau 包络给出著名的 Huber 型平滑：
\[
e_\lambda f(x)=
\begin{cases}
\dfrac{x^2}{2\lambda}, & |x|\le \lambda,\\[4pt]
|x|-\dfrac{\lambda}{2}, & |x|>\lambda.
\end{cases}
\]
把这个例子放在这里，是为了给 Boyd 读者一个最熟悉的有限维坍缩对象：不可微绝对值经过下确界卷积后变成了分段光滑的 Huber 损失。
\end{example}

\begin{example}[正则化分裂模型]
对
\[
F(x)=\varphi(Ax)+\psi(x),
\]
第 \ref{sec:5-4} 节中的次微分演算告诉我们如何写最优性条件，而本节则说明若把某一项替换为 Moreau 包络，就能得到更平滑的近似模型。把这个例子放在这里，是为了直接接通第 9 章 Moreau--Yosida 正则化与第 10 章变量分裂算法。
\end{example}

\subsection*{有限维对照}

Boyd 书里对 Huber 损失、平滑近似与变量分裂的讨论，经常以具体矩阵模型和一维图像展开。本节说明，这些对象背后的统一代数恰恰是定义~\ref{def:infimal-convolution} 与定理~\ref{thm:conjugate-infimal-convolution}。有限维里看上去是“技巧”的平滑与分裂，在无穷维里其实是一种由共轭理论控制的结构性变换。

\subsection*{习题（待补）}

\begin{exercise}[下确界卷积占位]\label{exr:infimal-convolution-placeholder}
补一题：对两个具体的凸函数计算其下确界卷积，并验证定理~\ref{thm:conjugate-infimal-convolution} 在该例中的公式。
\end{exercise}

\begin{exercise}[Moreau 包络桥接题占位]\label{exr:moreau-envelope-placeholder}
补一题：从例中的绝对值函数出发计算 Moreau 包络，并解释它为何会自然导向 Huber 平滑和近端算法。
\end{exercise}


\section{Fenchel-Rockafellar 对偶框架}

\label{sec:6-4}

前面三节把共轭、双共轭与下确界卷积准备齐全之后，对偶理论终于可以被写成一个真正可复用的模板：给定原问题，如何构造其对偶问题？为何对偶目标总是原问题的下界？何时下界可以提升为精确的最优值表达？这些问题在有限维教材中常以拉格朗日函数与 KKT 条件的形式出现，而在无穷维里，它们的统一上位版本就是 Fenchel--Rockafellar 对偶。第 \ref{sec:5-3} 节定义~\ref{def:subdifferential} 给出的局部最优性语言，以及第 \ref{sec:5-4} 节定理~\ref{thm:moreau-rockafellar-sum} 中已经出现的资格条件机制，都会在这里被收束为一个原--对偶总模板。第 \ref{sec:7-1} 节开始的广义不等式与第 \ref{sec:7-4} 节的无穷维 KKT，都会直接建立在这一节的框架之上。

\subsection{原问题、对偶问题与弱对偶}

\begin{definition}[Fenchel 型原--对偶对]\label{def:fenchel-primal-dual-pair}
设 $X,Y$ 为 Banach 空间，$A\colon X\to Y$ 为有界线性算子，$f\colon X\to \mathbb R\cup\{+\infty\}$ 与 $g\colon Y\to \mathbb R\cup\{+\infty\}$ 为 proper 凸泛函。定义\emph{原问题}
\[
p^\star:=\inf_{x\in X}\{f(x)+g(Ax)\},
\]
以及对应的\emph{对偶问题}
\[
d^\star:=\sup_{y^\ast\in Y^\ast}\{-f^\ast(-A^\ast y^\ast)-g^\ast(y^\ast)\}.
\]
数量
\[
p^\star-d^\star
\]
称为\emph{对偶间隙}。
\end{definition}

\begin{proposition}[弱对偶]\label{prop:fenchel-weak-duality}
对任意 $x\in X$ 与任意 $y^\ast\in Y^\ast$，都有
\[
-f^\ast(-A^\ast y^\ast)-g^\ast(y^\ast)\le f(x)+g(Ax).
\]
因此
\[
d^\star\le p^\star.
\]
\end{proposition}

\begin{proof}
对 $f$ 应用命题~\ref{prop:fenchel-young-inequality}，并令对偶变量取为 $-A^\ast y^\ast$，得
\[
f(x)+f^\ast(-A^\ast y^\ast)\ge \langle -A^\ast y^\ast,x\rangle=-\langle y^\ast,Ax\rangle.
\]
对 $g$ 应用同一不等式，得
\[
g(Ax)+g^\ast(y^\ast)\ge \langle y^\ast,Ax\rangle.
\]
两式相加即得
\[
f(x)+g(Ax)+f^\ast(-A^\ast y^\ast)+g^\ast(y^\ast)\ge 0,
\]
移项得到
\[
-f^\ast(-A^\ast y^\ast)-g^\ast(y^\ast)\le f(x)+g(Ax).
\]
再对左边取上确界、对右边取下确界，便得到
\[
d^\star\le p^\star.
\]
\end{proof}

\begin{remark}
命题~\ref{prop:fenchel-weak-duality} 的意义在于：对偶问题从来不会给出“过高的承诺”。它所产生的任何值，都是原问题最优值的合法下界。真正困难的是判断这个下界是否精确，也即对偶间隙是否为零；这正是资格条件要解决的问题。
\end{remark}

\subsection{资格条件}

\begin{definition}[Fenchel 对偶的资格条件]\label{def:constraint-qualification-fenchel}
设 $(f,g,A)$ 如定义~\ref{def:fenchel-primal-dual-pair}。若存在某个 $x_0\in \operatorname{dom}f$ 使得 $Ax_0\in \operatorname{dom}g$ 且 $g$ 在 $Ax_0$ 处连续，则称该三元组满足\emph{Fenchel 型资格条件}。
\end{definition}

\begin{remark}
定义~\ref{def:constraint-qualification-fenchel} 是本章采用的 Banach/Hilbert 常用版本。更一般的局部凸空间理论常把资格条件写成
\[
0\in \operatorname{core}(\operatorname{dom}g-A(\operatorname{dom}f)),
\]
这里的 core 正是第 \ref{sec:4-2} 节引入的代数内部。为了保持正文主线紧凑，本节把证明建立在连续性资格条件上，而把更一般版本留作前瞻备注。
\end{remark}

\subsection{Fenchel--Rockafellar 对偶定理}

\begin{theorem}[Fenchel--Rockafellar 对偶定理]\label{thm:fenchel-rockafellar-duality}
设 $X,Y$ 为 Banach 空间，$A\colon X\to Y$ 为有界线性算子，$f\colon X\to \mathbb R\cup\{+\infty\}$ 与 $g\colon Y\to \mathbb R\cup\{+\infty\}$ 为 proper、凸且下半连续的泛函。若定义~\ref{def:constraint-qualification-fenchel} 的资格条件成立，则
\[
\inf_{x\in X}\{f(x)+g(Ax)\}
=
\sup_{y^\ast\in Y^\ast}\{-f^\ast(-A^\ast y^\ast)-g^\ast(y^\ast)\}.
\]
换言之，在该资格条件下对偶间隙为零，即发生强对偶。
\end{theorem}

\begin{proof}
定义扰动函数
\[
F(x,y):=f(x)+g(Ax+y),\qquad (x,y)\in X\times Y,
\]
以及值函数
\[
\varphi(y):=\inf_{x\in X}F(x,y)=\inf_{x\in X}\{f(x)+g(Ax+y)\}.
\]
显然
\[
\varphi(0)=\inf_{x\in X}\{f(x)+g(Ax)\}=p^\star.
\]

首先说明 $\varphi$ 是 proper、凸且下半连续的，并且在 $0$ 处连续。proper 性来自资格条件中的点 $x_0$：由于 $f(x_0)<+\infty$ 且 $g$ 在 $Ax_0$ 处连续，故存在 $0$ 的邻域 $U\subset Y$ 使得对所有 $y\in U$，
\[
g(Ax_0+y)<+\infty.
\]
于是对所有 $y\in U$，
\[
\varphi(y)\le f(x_0)+g(Ax_0+y)<+\infty.
\]
又因 $g$ 在 $Ax_0$ 处连续，右侧随 $y\to 0$ 有界，从而 $\varphi$ 在 $0$ 的邻域上有上界。另一方面，作为一族仿射变换下凸函数的下确界，$\varphi$ 自身仍为凸。凸函数在内点附近若局部有上界，则连续，故 $\varphi$ 在 $0$ 处连续。由连续性可知 $\varphi$ 在 $0$ 处当然下半连续。

接着计算 $\varphi^\ast$。对任意 $y^\ast\in Y^\ast$，
\[
\varphi^\ast(y^\ast)
=
\sup_{y\in Y}\{\langle y^\ast,y\rangle-\varphi(y)\}
=
\sup_{x\in X,\ y\in Y}\{\langle y^\ast,y\rangle-f(x)-g(Ax+y)\}.
\]
令
\[
z:=Ax+y,
\]
则 $y=z-Ax$，于是
\[
\varphi^\ast(y^\ast)
=
\sup_{x\in X,\ z\in Y}\{\langle y^\ast,z-Ax\rangle-f(x)-g(z)\}.
\]
整理得
\[
\varphi^\ast(y^\ast)
=
\sup_{x\in X}\{-\langle A^\ast y^\ast,x\rangle-f(x)\}
\;+\;
\sup_{z\in Y}\{\langle y^\ast,z\rangle-g(z)\}.
\]
因此
\[
\varphi^\ast(y^\ast)=f^\ast(-A^\ast y^\ast)+g^\ast(y^\ast).
\]

由于 $\varphi$ 是 proper、凸且在 $0$ 处连续，故由第 \ref{sec:6-2} 节定理~\ref{thm:fenchel-moreau}，
\[
\varphi(0)=\varphi^{\ast\ast}(0)
=
\sup_{y^\ast\in Y^\ast}\{-\varphi^\ast(y^\ast)\}.
\]
代入上一步的表达式，得到
\[
p^\star
=
\sup_{y^\ast\in Y^\ast}\{-f^\ast(-A^\ast y^\ast)-g^\ast(y^\ast)\}
=
d^\star.
\]
故强对偶成立。
\end{proof}

\begin{remark}
定理~\ref{thm:fenchel-rockafellar-duality} 的证明清楚地表明：资格条件的作用，是保证扰动值函数 $\varphi$ 在原点附近足够正则，从而能够通过定理~\ref{thm:fenchel-moreau} 被双共轭完全重建。也就是说，强对偶并不是神奇的值相等，而是“扰动问题在原点没有病态缺口”这一事实的结果。
\end{remark}

\subsection{最优性条件与 KKT 接口}

\begin{corollary}[Fenchel 对偶的次微分最优性条件]\label{cor:fenchel-kkt-subdifferential}
在定理~\ref{thm:fenchel-rockafellar-duality} 的假设下，设 $\bar x\in X$ 与 $\bar y^\ast\in Y^\ast$ 分别是原问题与对偶问题的最优解。则下列条件等价于原--对偶最优性：
\[
-A^\ast \bar y^\ast\in \partial f(\bar x),
\qquad
\bar y^\ast\in \partial g(A\bar x).
\]
等价地，
\[
0\in \partial f(\bar x)+A^\ast \partial g(A\bar x).
\]
\end{corollary}

\begin{proof}
由最优性与强对偶，
\[
f(\bar x)+g(A\bar x)
=
-f^\ast(-A^\ast \bar y^\ast)-g^\ast(\bar y^\ast).
\]
整理为
\[
f(\bar x)+f^\ast(-A^\ast \bar y^\ast)
=
-\langle \bar y^\ast,A\bar x\rangle,
\qquad
g(A\bar x)+g^\ast(\bar y^\ast)
=
\langle \bar y^\ast,A\bar x\rangle.
\]
于是对 $f$ 与 $g$ 分别都在 Fenchel--Young 不等式中取等。由命题~\ref{prop:fenchel-young-equality-subgradient}，便得到
\[
-A^\ast \bar y^\ast\in \partial f(\bar x),
\qquad
\bar y^\ast\in \partial g(A\bar x).
\]
反过来，若这两条次微分条件成立，则再次应用命题~\ref{prop:fenchel-young-equality-subgradient} 可知两边分别在 Fenchel--Young 不等式中取等，故原目标值与对偶目标值相等。结合命题~\ref{prop:fenchel-weak-duality}，便知两者均为最优。
\end{proof}

\begin{remark}
推论~\ref{cor:fenchel-kkt-subdifferential} 是第 \ref{sec:7-1} 节广义不等式和第 \ref{sec:7-4} 节无穷维 KKT 系统的直接前导接口。有限维中常见的拉格朗日乘子、驻点条件与互补松弛，在本书后文都会被重新写回这里的次微分与伴随算子语言。
\end{remark}

\subsection{典型例子}

\begin{example}[线性算子复合的原--对偶模板]
设 $A\colon X\to Y$ 为有界线性算子，原问题
\[
\inf_{x\in X}\{f(x)+g(Ax)\}
\]
的对偶正是
\[
\sup_{y^\ast\in Y^\ast}\{-f^\ast(-A^\ast y^\ast)-g^\ast(y^\ast)\}.
\]
把这个例子放在这里，是为了强调第 \ref{sec:7-2} 节抽象标准型、第 \ref{sec:10-3} 节 ADMM 以及大量逆问题模型，其实都只是这一统一模板的具体化。
\end{example}

\begin{example}[有限维拉格朗日对偶]
当 $X=\mathbb R^n$、$Y=\mathbb R^m$ 且 $A$ 是矩阵时，定义~\ref{def:fenchel-primal-dual-pair} 退化为 Boyd 书中最熟悉的有限维对偶模板。把这个例子放在这里，是为了说明 Boyd 第 5 章的拉格朗日对偶并不是另一套平行理论，而是 Fenchel--Rockafellar 对偶在有限维 Hilbert 空间中的坍缩。
\end{example}

\begin{example}[资源约束的价差解释]
在管理科学模型里，$g(Ax)$ 常用来编码资源或风险约束。对偶变量 $y^\ast$ 则可解释为资源影子价格或边际罚率。把这个例子放在这里，是为了强调对偶变量不只是证明中的中介量；在许多应用场景中，它本身就携带了可解释的经济或工程信息。
\end{example}

\subsection*{有限维对照}

Boyd 书中关于强对偶、对偶间隙与 KKT 的叙述，往往借助 Slater 条件和拉格朗日函数在 $\mathbb R^n$ 中展开。本节说明，这些有限维公式真正依赖的核心结构是：定义~\ref{def:fenchel-primal-dual-pair} 给出的原--对偶模板、命题~\ref{prop:fenchel-weak-duality} 给出的下界机制，以及定理~\ref{thm:fenchel-rockafellar-duality} 所控制的资格条件。有限维之所以显得简单，是因为很多正则性在那里自动满足，而不是因为对偶理论只属于矩阵世界。

\subsection*{习题（待补）}

\begin{exercise}[弱对偶占位]\label{exr:weak-duality-placeholder}
补一题：对一个具体的原--对偶对，逐步验证命题~\ref{prop:fenchel-weak-duality}，并写出它给出的显式对偶下界。
\end{exercise}

\begin{exercise}[Fenchel--Rockafellar 桥接题占位]\label{exr:fenchel-rockafellar-placeholder}
补一题：在有限维矩阵模型中应用定理~\ref{thm:fenchel-rockafellar-duality}，并把推论~\ref{cor:fenchel-kkt-subdifferential} 翻译成 Boyd 风格的 KKT 条件。
\end{exercise}

